\chapter{Decisioni bayesiane e classificazione binaria}

L'idea principale della statistica inferenziale è che, osservato un campione sufficientemente ampio di una popolazione,
si possono dedurre alcune caratteristiche che devono essere rispettate da qualsiasi altro campione. Spesso si ha anche
una certa conoscenza a priori: per esempio si sa che il campione proviene da una popolazione la cui distribuzione è nota
a meno di un insieme di parametri, cioè da una \textbf{famiglia di distribuzioni} indicizzate da un parametro $\theta$.
La domanda diventa: come si elaborano i dati per dedurre il valore del parametro?

Se il parametro può assumere solo un numero finito di valori, la domanda è una \textbf{decisione}: quale fra $M$
spiegazioni dei dati è quella vera? È il problema di un ricevitore che deve capire quale simbolo è stato trasmesso, di
un filtro che deve classificare un messaggio come posta indesiderata o no, di un test medico. In questo capitolo lo
risolviamo nel caso \textbf{bayesiano}, in cui si conoscono le probabilità a priori delle varie spiegazioni.

% ============================================================
\section{Regole di decisione e rischio bayesiano}
% ============================================================

Si abbia un insieme di dati $\vett{x} \in \mathcal{X}^n$, realizzazione di un vettore aleatorio $\vett{X}$. I dati
possono seguire una fra $M$ diverse leggi di probabilità: c'è un insieme di $M$ \textbf{ipotesi} mutuamente esclusive
$\{H_i\}_{i=1}^{M}$, ciascuna delle quali definisce una diversa legge condizionata per i dati,
\[ p_{\vett{X}}(\vett{x} \mid H_i) , \qquad i = 1, 2, \ldots, M . \]
In altre parole $\vett{X}$ ha una distribuzione che appartiene a una famiglia di pmf $p_{\vett{X}}(\vett{x} \mid
\theta)$, dove il valore del parametro $\theta$ (qui, l'indice $i$ dell'ipotesi) non è noto. Si suppone inoltre di
conoscere le \textbf{probabilità a priori} $\{P(H_i)\}_{i=1}^{M}$ di questi «stati della natura».

\dfn{Regola di decisione}{
  Una \term{regola di decisione} è una corrispondenza
  \[ D : \vett{x} \in \mathcal{X}^n \;\longmapsto\; D(\vett{x}) \in \{1, 2, \ldots, M\} \]
  che, osservati i dati, decide quale dei possibili stati della natura è quello effettivamente attivo.
}

Dare una regola di decisione equivale a dividere l'insieme di tutti i dati possibili in $M$ \term{regioni di
decisione} $\Omega_i = \{\vett{x} : D(\vett{x}) = i\}$: se i dati cadono in $\Omega_i$ si decide per l'ipotesi $H_i$.
Le regioni sono una partizione di $\mathcal{X}^n$ (Capitolo~1): non si sovrappongono e lo ricoprono tutto, perché per
ogni dato si prende una e una sola decisione.

\begin{figure}[H]
\centering
\begin{tikzpicture}
  \fill[mspiA!12] (0,0) -- (3.4,0) .. controls (3.0,1.2) and (4.6,2.0) .. (4.0,3.2) -- (0,3.2) -- cycle;
  \fill[mspiB!14] (3.4,0) .. controls (3.0,1.2) and (4.6,2.0) .. (4.0,3.2) -- (7.6,3.2) -- (7.6,0) -- cycle;
  \draw[mspigray, line width=0.6pt] (0,0) rectangle (7.6,3.2);
  \draw[black, line width=1pt] (3.4,0) .. controls (3.0,1.2) and (4.6,2.0) .. (4.0,3.2);
  \node[font=\small] at (1.7,2.5) {$\Omega_1$: si decide $H_1$};
  \node[font=\small] at (6.0,2.5) {$\Omega_2$: si decide $H_2$};
  \foreach \p in {(0.8,0.9), (1.5,1.6), (2.4,0.6), (2.9,2.0), (1.2,2.0)} \fill[mspiA] \p circle (1.6pt);
  \foreach \p in {(5.0,0.8), (6.3,1.5), (4.7,1.9), (6.8,0.5), (5.6,2.0)} \fill[mspiB] \p circle (1.6pt);
  \node[font=\scriptsize, mspigray] at (7.25,0.25) {$\mathcal{X}^n$};
  \node[font=\small, anchor=west, text width=6.4cm] at (8.2,1.6) {L'insieme $\mathcal{X}^n$ di tutti i dati possibili
    è diviso in due regioni. Progettare la regola di decisione vuol dire scegliere da che parte mettere ogni
    possibile $\vett{x}$, cioè dove tracciare il confine.};
\end{tikzpicture}
\caption{Una regola di decisione fra due ipotesi è una partizione dei dati in due regioni.}
\end{figure}

Per scegliere la regola migliore bisogna dire quanto costano gli errori. Non tutti gli errori sono uguali: scambiare
un messaggio importante per posta indesiderata è peggio del contrario.

\dfn{Matrice dei costi}{
  La \term{matrice dei costi} è la matrice $M \times M$
  \[ \matr{C} = \begin{bmatrix} C_{11} & C_{12} & \cdots & C_{1M} \\ C_{21} & C_{22} & \cdots & C_{2M} \\ \vdots &
     \vdots & \ddots & \vdots \\ C_{M1} & C_{M2} & \cdots & C_{MM} \end{bmatrix} , \]
  dove $C_{ij}$ è il costo associato all'evento «si prende la decisione per l'ipotesi $i$, $D(\vett{x}) = i$, e lo
  stato della natura è $H_j$», cioè l'ipotesi vera è la $j$.
}

\dfn{Rischio medio bayesiano}{
  Il \term{rischio medio bayesiano} di una regola di decisione $D$ è il costo medio:
  \[ R = \sum_{i=1}^{M} \sum_{j=1}^{M} C_{ij}\, P\{D(\vett{X}) = i,\ H = H_j\} . \]
  Data una matrice dei costi $\matr{C}$, una \term{regola di decisione ottima} è una corrispondenza $D(\vett{x})$ che
  minimizza il rischio bayesiano.
}

Il rischio è la media di una variabile aleatoria, il costo: il costo $C_{ij}$ si paga con la probabilità che si decida
$i$ quando è vera $j$. Il caso più comune è quello in cui le decisioni giuste non costano nulla e quelle sbagliate
costano tutte allo stesso modo.

\teorema[rischioerrore]{Costi uniformi: il rischio è la probabilità di errore}{
  Se $C_{ii} = 0$ per $i = 1, 2, \ldots, M$ e $C_{ij} = 1$ per ogni $i \neq j$, allora il rischio medio bayesiano
  coincide con la probabilità di commettere un errore nella classificazione:
  \[ R = \sum_{i=1}^{M} \sum_{\substack{j=1 \\ j \neq i}}^{M} P\{D(\vett{X}) = i,\ H = H_j\} = P(e) . \]
}

\begin{dimostrazione}
\passo{} I termini con $i = j$ hanno costo zero e spariscono; gli altri hanno costo 1 e restano con la loro
probabilità. Gli eventi $\{D(\vett{X}) = i, H = H_j\}$ con $i \neq j$ sono incompatibili (in ogni prova si prende una
sola decisione ed è vera una sola ipotesi), e la loro unione è l'evento «la decisione è sbagliata»: per l'additività la
somma delle loro probabilità è $P(e)$.
\end{dimostrazione}

% ============================================================
\section{La classificazione binaria a minima probabilità di errore}
% ============================================================

Supponiamo ora $M = 2$, $C_{11} = C_{22} = 0$ e $C_{12} = C_{21} = 1$. Per il teorema precedente
\[ R = P\{D(\vett{X}) = 2,\ H_1\} + P\{D(\vett{X}) = 1,\ H_2\} = P(e) . \]
Progettare una regola di decisione vuol dire determinare una partizione di $\mathcal{X}^n$ in due sottoinsiemi
$\Omega_1$ e $\Omega_2$, con
\[ D(\vett{x}) = \begin{cases} 1 & \text{se } \vett{x} \in \Omega_1 , \\ 2 & \text{se } \vett{x} \in \Omega_2 ,
   \end{cases} \qquad\text{e quindi}\qquad P(e) = P\{\vett{X} \in \Omega_1,\ H_2\} + P\{\vett{X} \in \Omega_2,\ H_1\} . \]
Cerchiamo la regola ottima, cioè quella con la minima probabilità di errore. Supponiamo prima che le osservazioni siano
un vettore aleatorio \textbf{discreto}, con pmf condizionate $p_{\vett{X}}(\vett{x} \mid H_i)$ note.

\teorema[mapbinario]{Regola a minima probabilità di errore}{
  La regola di decisione che minimizza la probabilità di errore è
  \[ \vett{x} \in \Omega_1 \iff p_{\vett{X}}(\vett{x} \mid H_1)\, P(H_1) > p_{\vett{X}}(\vett{x} \mid H_2)\, P(H_2) , \]
  o, in modo equivalente,
  \[ L(\vett{x}) = \frac{p_{\vett{X}}(\vett{x} \mid H_1)}{p_{\vett{X}}(\vett{x} \mid H_2)} \decidi{H_1}{H_2}
     \frac{P(H_2)}{P(H_1)} = \eta , \]
  dove il simbolo $\decidi{H_1}{H_2}$ si legge: «se vale $>$ si decide $H_1$, se vale $<$ si decide $H_2$». La
  quantità $L(\vett{x})$ si chiama \term{rapporto di verosimiglianza} (\emph{likelihood ratio}) fra le due ipotesi.
}

\begin{dimostrazione}[la probabilità di errore è il complemento di quella di decisione corretta]
\passo{Si scrivono le due probabilità congiunte} Per la probabilità composta,
\[ P\{\vett{X} \in \Omega_1,\ H_2\} = \sum_{\vett{x} \in \Omega_1} P\{\vett{X} = \vett{x} \mid H_2\}\, P(H_2) , \qquad
   P\{\vett{X} \in \Omega_2,\ H_1\} = \sum_{\vett{x} \in \Omega_2} P\{\vett{X} = \vett{x} \mid H_1\}\, P(H_1) . \]
\passo{Si passa al complemento} Le due regioni si spartiscono tutti i dati, e la somma di una pmf su tutti i dati vale
1: quindi $\sum_{\vett{x} \in \Omega_1} p_{\vett{X}}(\vett{x} \mid H_2) = 1 - \sum_{\vett{x} \in \Omega_2}
p_{\vett{X}}(\vett{x} \mid H_2)$, e lo stesso per $H_1$. Sostituendo e usando $P(H_1) + P(H_2) = 1$:
\[ P(e) = 1 - \left[ \sum_{\vett{x} \in \Omega_1} p_{\vett{X}}(\vett{x} \mid H_1)\, P(H_1) + \sum_{\vett{x} \in
   \Omega_2} p_{\vett{X}}(\vett{x} \mid H_2)\, P(H_2) \right] . \]
La quantità fra parentesi quadre è la probabilità di una decisione corretta.
\passo{Si massimizza un dato alla volta} $P(e)$ è minima quando la parentesi quadra è massima. Ogni vettore $\vett{x}$
va messo o in $\Omega_1$ o in $\Omega_2$, non in entrambe: se va in $\Omega_1$ contribuisce alla parentesi con
$p_{\vett{X}}(\vett{x} \mid H_1) P(H_1)$, se va in $\Omega_2$ con $p_{\vett{X}}(\vett{x} \mid H_2) P(H_2)$. La scelta
migliore è metterlo dove il suo contributo è più grande: è la regola dell'enunciato.
\passo{Il rapporto} Dividendo i due membri della disuguaglianza per le quantità positive $p_{\vett{X}}(\vett{x} \mid
H_2)$ e $P(H_1)$ si ottiene la forma con il rapporto di verosimiglianza. (Se i due membri sono uguali, le due decisioni
danno lo stesso contributo e si può scegliere a piacere.)
\end{dimostrazione}

La stessa regola si legge anche in un altro modo, che le dà il nome. Per la legge di Bayes
\[ P\{H_i \mid \vett{X} = \vett{x}\} = \frac{P\{\vett{X} = \vett{x} \mid H_i\}\, P(H_i)}{P\{\vett{X} = \vett{x}\}} =
   \frac{p_{\vett{X}}(\vett{x} \mid H_i)\, P(H_i)}{p_{\vett{X}}(\vett{x})} , \]
e il denominatore è lo stesso per le due ipotesi: confrontare $p_{\vett{X}}(\vett{x} \mid H_i) P(H_i)$ equivale a
confrontare le probabilità a posteriori.

\dfn{Regola MAP e regola ML}{
  La regola a minima probabilità di errore si chiama anche \term{regola MAP} (\emph{Massima probabilità A
  Posteriori}): decide per l'ipotesi la cui probabilità a posteriori, dati i dati osservati, è massima. Se le due
  ipotesi sono a priori equiprobabili, la soglia vale $\eta = 1$ e la regola diventa la \term{regola a massima
  verosimiglianza}, o \term{regola ML} (\emph{Maximum Likelihood}): decide per l'ipotesi che rende più probabili i dati
  osservati.
}

Le probabilità a priori spostano la soglia: se $H_2$ è a priori molto più probabile di $H_1$, la soglia $\eta$ è grande,
e per decidere $H_1$ servono dati che la favoriscano nettamente. Le \textbf{probabilità di errore condizionate} sono
\[ P(e \mid H_1) = P\{L(\vett{X}) < \eta \mid H_1\} , \qquad P(e \mid H_2) = P\{L(\vett{X}) > \eta \mid H_2\} , \]
e per la legge della probabilità totale
\[ P(e) = P(H_1)\, P(e \mid H_1) + P(H_2)\, P(e \mid H_2) . \]

% ============================================================
\section{Esempio: da quale sorgente vengono i bit?}
% ============================================================

Le osservazioni sono $n$ variabili binarie i.i.d., estratte con la stessa probabilità da una sorgente con
$P\{X_i = 1\} = p_1$ (ipotesi $H_1$) o da una sorgente con $P\{X_i = 1\} = p_2$ (ipotesi $H_2$), con $p_1 > p_2$.
Una sequenza precisa con $w$ uni e $n - w$ zeri ha probabilità $p_i^{w} (1 - p_i)^{n - w}$ (i bit sono indipendenti),
quindi
\[ p_{\vett{X}}(\vett{x} \mid H_i) = p_i^{w_H(\vett{x})}\, (1 - p_i)^{n - w_H(\vett{x})} , \]
dove $w_H(\vett{x})$ è il \term{peso di Hamming} della sequenza binaria osservata, cioè il numero dei suoi «1».

\teorema[testhamming]{Il test sul peso di Hamming}{
  Con ipotesi equiprobabili, il test a minima probabilità di errore fra le due sorgenti è
  \[ w_H(\vett{x}) \decidi{H_1}{H_2} \eta_1 = n\, \frac{\ln \dfrac{1 - p_2}{1 - p_1}}{\ln \dfrac{p_1 (1 - p_2)}{p_2 (1 -
     p_1)}} : \]
  si decide per la sorgente con più uni se la sequenza ha abbastanza uni.
}

\begin{dimostrazione}
\passo{Il rapporto di verosimiglianza} Con $P(H_1) = P(H_2) = \frac12$ la soglia è $\eta = 1$, e raccogliendo le
potenze con lo stesso esponente:
\[ L(\vett{x}) = \frac{p_1^{w_H(\vett{x})} (1 - p_1)^{n - w_H(\vett{x})}}{p_2^{w_H(\vett{x})} (1 - p_2)^{n -
   w_H(\vett{x})}} = \left( \frac{p_1}{p_2} \right)^{w_H(\vett{x})} \left( \frac{1 - p_1}{1 - p_2} \right)^{n -
   w_H(\vett{x})} \decidi{H_1}{H_2} 1 . \]
\passo{Il logaritmo} Il logaritmo è una funzione crescente: applicato a entrambi i membri non cambia il verso della
disuguaglianza, e trasforma le potenze in prodotti ($\ln a^b = b \ln a$) e il prodotto in somma. Con $w = w_H(\vett{x})$:
\[ w \ln \frac{p_1}{p_2} + (n - w) \ln \frac{1 - p_1}{1 - p_2} \decidi{H_1}{H_2} 0 \quad\iff\quad w \left[ \ln
   \frac{p_1}{p_2} - \ln \frac{1 - p_1}{1 - p_2} \right] \decidi{H_1}{H_2} n \ln \frac{1 - p_2}{1 - p_1} , \]
dove si è usato $-\ln \frac{1 - p_1}{1 - p_2} = \ln \frac{1 - p_2}{1 - p_1}$.
\passo{Si isola il peso} La differenza di logaritmi è il logaritmo del quoziente:
$w \ln \frac{p_1 (1 - p_2)}{p_2 (1 - p_1)} \decidi{H_1}{H_2} n \ln \frac{1 - p_2}{1 - p_1}$. Siccome $p_1 > p_2$, si ha
$\frac{p_1}{p_2} > 1$ e $\frac{1 - p_2}{1 - p_1} > 1$: tutti i logaritmi sono positivi, e dividendo per il
coefficiente di $w$ il verso non cambia.
\end{dimostrazione}

Si noti una cosa importante: per decidere non serve conoscere tutta la sequenza, ma solo quanti uni contiene. Una
quantità calcolata dai dati che contiene tutto ciò che serve per la decisione si chiama, nel gergo della statistica
inferenziale, \term{statistica sufficiente}: qui è $w_H(\vett{x})$.

Per le prestazioni, si osservi che sotto l'ipotesi $H_i$ il peso $w_H(\vett{X})$ conta i successi in $n$ prove
indipendenti con probabilità $p_i$: è \textbf{binomiale} $\mathcal{B}(n, p_i)$ (Capitolo~5). Se $\eta_1$ non è intero,
\[ P(e \mid H_1) = P\{w_H(\vett{X}) < \eta_1 \mid H_1\} = \sum_{i=0}^{\lfloor \eta_1 \rfloor} \binom{n}{i} p_1^i (1 -
   p_1)^{n - i} , \qquad P(e \mid H_2) = P\{w_H(\vett{X}) > \eta_1 \mid H_2\} = \sum_{i = \lfloor \eta_1 \rfloor +
   1}^{n} \binom{n}{i} p_2^i (1 - p_2)^{n - i} , \]
dove $\lfloor \eta_1 \rfloor$ è la \term{parte intera} di $\eta_1$ (il più grande intero che non la supera), e
\[ P(e) = P(e \mid H_1)\, P(H_1) + P(e \mid H_2)\, P(H_2) = \frac12 P(e \mid H_1) + \frac12 P(e \mid H_2) . \]

\begin{esempio}{Dieci bit, sorgenti con $p_1 = 0{,}6$ e $p_2 = 0{,}2$}
La soglia è
\[ \eta_1 = 10\, \frac{\ln (0{,}8/0{,}4)}{\ln \frac{0{,}6 \cdot 0{,}8}{0{,}2 \cdot 0{,}4}} = 10\, \frac{\ln 2}{\ln 6} =
   10 \cdot \frac{0{,}693}{1{,}792} \approx 3{,}87 : \]
si decide $H_1$ se la sequenza contiene almeno 4 uni. Le probabilità di errore sono $P(e \mid H_1) = P\{w \leq 3 \mid
p = 0{,}6\} \approx 0{,}055$ e $P(e \mid H_2) = P\{w \geq 4 \mid p = 0{,}2\} \approx 0{,}121$, quindi $P(e) \approx
0{,}088$. Nella figura, per ogni peso $w$, le due barre sono $\frac12 P(w \mid H_1)$ e $\frac12 P(w \mid H_2)$: la
regola ottima sceglie, per ogni $w$, la barra più alta, e l'errore è la somma delle barre più basse.
\begin{center}
\begin{tikzpicture}
\begin{axis}[width=14cm, height=5.6cm, ybar, bar width=8pt, xmin=-0.7, xmax=10.7, ymin=0, ymax=0.175,
  xtick={0,1,...,10}, ytick={0.05,0.1,0.15}, yticklabels={$0{,}05$,$0{,}1$,$0{,}15$}, xlabel={peso di Hamming $w$},
  axis lines=left, tick label style={font=\scriptsize}, label style={font=\small}, grid=major, grid style={mspigray!20},
  legend style={font=\scriptsize, at={(0.99,0.98)}, anchor=north east}, clip=false]
  \addplot[fill=mspiA!60, draw=mspiA!80!black] coordinates {(0,0.0001) (1,0.0008) (2,0.0053) (3,0.0212) (4,0.0557)
    (5,0.1003) (6,0.1254) (7,0.1075) (8,0.0605) (9,0.0202) (10,0.0030)};
  \addplot[fill=mspiB!60, draw=mspiB!80!black] coordinates {(0,0.0537) (1,0.1342) (2,0.1510) (3,0.1007) (4,0.0440)
    (5,0.0132) (6,0.0028) (7,0.0004) (8,0.0000) (9,0.0000) (10,0.0000)};
  \draw[black, dashed, line width=1pt] (axis cs:3.87,0) -- (axis cs:3.87,0.165);
  \node[font=\small, anchor=east] at (axis cs:3.7,0.165) {si decide $H_2$};
  \node[font=\small, anchor=west] at (axis cs:4.05,0.165) {si decide $H_1$};
  \legend{{$\frac12 P(w \mid H_1)$}, {$\frac12 P(w \mid H_2)$}}
\end{axis}
\end{tikzpicture}
\end{center}
Che la soglia 4 sia la migliore si può verificare provando tutte le soglie possibili («si decide $H_1$ se $w \geq t$»):
\begin{center}
\renewcommand{\arraystretch}{1.25}
\begin{tabular}{@{}c|ccccccc@{}}
  soglia $t$ & 1 & 2 & 3 & \textbf{4} & 5 & 6 & 7 \\ \hline
  $P(e)$ & 0,446 & 0,313 & 0,167 & \textbf{0,088} & 0,100 & 0,187 & 0,309
\end{tabular}
\end{center}
\end{esempio}

% ============================================================
\section{Dati continui}
% ============================================================

Se i dati sono estratti da una fra $M$ leggi continue, con densità condizionate $\{f_{\vett{X} \mid H_i}(\vett{x} \mid
H_i)\}_{i=1}^{M}$, l'unica differenza è che le somme sulle regioni diventano integrali; per esempio, con $M = 2$,
\[ P\{\vett{X} \in \Omega_1 \mid H_1\} = \int_{\Omega_1} f_{\vett{X} \mid H_1}(\vett{x} \mid H_1) \dd\vett{x} , \qquad
   P\{\vett{X} \in \Omega_2 \mid H_2\} = \int_{\Omega_2} f_{\vett{X} \mid H_2}(\vett{x} \mid H_2) \dd\vett{x} , \]
dove $\dd\vett{x}$ indica un integrale multiplo, su tutte le $n$ componenti.

\teorema[mapcontinuo]{Regola a minima probabilità di errore per dati continui}{
  \[ \vett{x} \in \Omega_1 \iff f_{\vett{X} \mid H_1}(\vett{x} \mid H_1)\, P(H_1) > f_{\vett{X} \mid H_2}(\vett{x} \mid
     H_2)\, P(H_2) , \qquad\text{cioè}\qquad L(\vett{x}) = \frac{f_{\vett{X} \mid H_1}(\vett{x} \mid H_1)}{f_{\vett{X}
     \mid H_2}(\vett{x} \mid H_2)} \decidi{H_1}{H_2} \frac{P(H_2)}{P(H_1)} . \]
  La quantità $L(\vett{x})$ è ancora il rapporto di verosimiglianza fra le due ipotesi.
}

\begin{dimostrazione}
\passo{} È lo stesso ragionamento del caso discreto, con gli integrali al posto delle somme:
\[ P(e) = 1 - \left[ \int_{\Omega_1} f_{\vett{X} \mid H_1}(\vett{x} \mid H_1)\, P(H_1) \dd\vett{x} + \int_{\Omega_2}
   f_{\vett{X} \mid H_2}(\vett{x} \mid H_2)\, P(H_2) \dd\vett{x} \right] , \]
e la parentesi è massima se ogni $\vett{x}$ viene messo nella regione in cui la funzione integranda è più grande.
\end{dimostrazione}

% ============================================================
\section{Esempio: la media di una popolazione gaussiana}
% ============================================================

I dati, con la stessa probabilità, sono la realizzazione di un vettore gaussiano con componenti indipendenti, tutte con
la stessa varianza $\sigma^2$, e medie $\mu_1$ (ipotesi $H_1$) oppure $\mu_2 < \mu_1$ (ipotesi $H_2$). Per esempio, $n$
misure rumorose di una tensione che può valere $\mu_1$ o $\mu_2$. La densità dei dati (Capitolo~16) è
\[ f_{\vett{X} \mid H_i}(\vett{x} \mid H_i) = \prod_{k=1}^{n} \frac{1}{\sqrt{2\pi\sigma^2}}\, e^{-\frac{(x_k -
   \mu_i)^2}{2\sigma^2}} . \]

\teorema[testmedia]{Il test sulla media}{
  Il test a minima probabilità di errore confronta la media campione con il punto di mezzo fra le due medie:
  \[ \frac{1}{n} \sum_{k=1}^{n} x_k \decidi{H_1}{H_2} \frac{\mu_1 + \mu_2}{2} , \]
  e la sua probabilità di errore è
  \[ P(e) = Q\left( \sqrt{n}\, \frac{\mu_1 - \mu_2}{2\sigma} \right) \;\xrightarrow{\;n \to \infty\;}\; 0 . \]
}

\begin{dimostrazione}
\passo{Il rapporto di verosimiglianza} Con ipotesi equiprobabili la soglia è 1. Il prodotto delle costanti si
semplifica, e il quoziente di esponenziali è l'esponenziale della differenza:
\[ L(\vett{x}) = \frac{f_{\vett{X} \mid H_1}(\vett{x} \mid H_1)}{f_{\vett{X} \mid H_2}(\vett{x} \mid H_2)} =
   e^{\sum_{k=1}^{n} \frac{(x_k - \mu_2)^2 - (x_k - \mu_1)^2}{2\sigma^2}} \decidi{H_1}{H_2} 1 . \]
\passo{Il logaritmo} Prendendo il logaritmo naturale e moltiplicando per $2\sigma^2 > 0$ (il verso non cambia):
\[ \sum_{k=1}^{n} \left[ (x_k - \mu_2)^2 - (x_k - \mu_1)^2 \right] \decidi{H_1}{H_2} 0 . \]
\passo{Si sviluppano i quadrati} I termini $x_k^2$ si cancellano:
\[ \sum_{k=1}^{n} \left[ x_k^2 + \mu_2^2 - 2 x_k \mu_2 - x_k^2 - \mu_1^2 + 2 x_k \mu_1 \right] = 2(\mu_1 - \mu_2)
   \sum_{k=1}^{n} x_k + n(\mu_2^2 - \mu_1^2) \decidi{H_1}{H_2} 0 . \]
\passo{Si isola la somma} Con $\mu_1^2 - \mu_2^2 = (\mu_1 + \mu_2)(\mu_1 - \mu_2)$:
$2(\mu_1 - \mu_2) \sum_k x_k \decidi{H_1}{H_2} n (\mu_1 + \mu_2)(\mu_1 - \mu_2)$. Siccome $\mu_1 - \mu_2 > 0$, si può
dividere per $2n(\mu_1 - \mu_2)$ senza cambiare il verso: è il test dell'enunciato.
\passo{La statistica del test} Sotto $H_i$, $Z_n = \frac{1}{n}\sum_k X_k$ è una combinazione lineare di gaussiane
indipendenti, quindi gaussiana (Capitolo~16), con media e varianza (Capitolo~18)
\[ E[Z_n \mid H_i] = \frac{1}{n} \sum_{k=1}^{n} \mu_i = \mu_i , \qquad \mathrm{VAR}[Z_n \mid H_i] = \frac{1}{n^2}
   \sum_{k=1}^{n} \sigma^2 = \frac{\sigma^2}{n} . \]
\passo{L'errore sotto $H_1$} Con la soglia $\eta = \frac{\mu_1 + \mu_2}{2}$, si standardizza e si usa la funzione $Q$
(Capitolo~15):
\[ P(e \mid H_1) = P\{Z_n < \eta \mid H_1\} = 1 - Q\left( \frac{\eta - \mu_1}{\sigma / \sqrt{n}} \right) = 1 - Q\left(
   \sqrt{n}\, \frac{\mu_2 - \mu_1}{2\sigma} \right) = Q\left( \sqrt{n}\, \frac{\mu_1 - \mu_2}{2\sigma} \right) , \]
perché $\eta - \mu_1 = \frac{\mu_2 - \mu_1}{2} < 0$ e $Q(-x) = 1 - Q(x)$.
\passo{L'errore sotto $H_2$}
\[ P(e \mid H_2) = P\{Z_n > \eta \mid H_2\} = Q\left( \frac{\eta - \mu_2}{\sigma / \sqrt{n}} \right) = Q\left( \sqrt{n}\,
   \frac{\mu_1 - \mu_2}{2\sigma} \right) . \]
\passo{L'errore totale} Le due probabilità condizionate sono uguali, quindi anche la loro media pesata, $P(e) =
P(H_1) P(e \mid H_1) + P(H_2) P(e \mid H_2)$, vale $Q\left( \sqrt{n} \frac{\mu_1 - \mu_2}{2\sigma} \right)$. Per $n \to
\infty$ l'argomento di $Q$ diverge, perché $\mu_1 - \mu_2 > 0$ e $\sigma > 0$, e $Q(+\infty) = 0$.
\end{dimostrazione}

Il risultato è molto intuitivo: si fa la media delle misure e si guarda se è più vicina a $\mu_1$ o a $\mu_2$. Anche qui
la statistica sufficiente è una sola, $\sum_k x_k$ (o la media campione): delle $n$ misure conta solo la somma. Le
slide indicano con $\eta_{Z_n}$ la media di $Z_n$; qui la chiamiamo $\mu_i$ per non confonderla con la soglia $\eta$.
Con $n = 1$, $\mu_1 = A$ e $\mu_2 = -A$ si ritrova la trasmissione binaria del Capitolo~15: soglia 0 e $P_e =
Q\left(\frac{A}{\sigma}\right)$. Ripetere la misura $n$ volte equivale a moltiplicare per $\sqrt{n}$ il rapporto fra la
distanza delle medie e la deviazione standard.

\begin{figure}[H]
\centering
\begin{tikzpicture}
\begin{groupplot}[group style={group size=2 by 1, horizontal sep=1.2cm}, width=8cm, height=5cm, quaderno,
  axis y line=none, xmin=-1.6, xmax=2.6, ymin=0, ymax=1.75, xtick={0,0.5,1}, xticklabels={$\mu_2$,$\eta$,$\mu_1$},
  xlabel=$z$, legend style={font=\scriptsize, at={(0.99,0.99)}, anchor=north east}]
\nextgroupplot[title={$n = 1$: $P(e) = Q(0{,}5) \approx 0{,}31$}]
  \addplot[fill=mspiB!35, draw=none, domain=0.5:2.6, samples=40, forget plot] {exp(-x^2/2)/sqrt(2*pi)} \closedcycle;
  \addplot[fill=mspiA!35, draw=none, domain=-1.6:0.5, samples=40, forget plot] {exp(-(x-1)^2/2)/sqrt(2*pi)} \closedcycle;
  \addplot[mspiA, domain=-1.6:2.6, samples=101] {exp(-(x-1)^2/2)/sqrt(2*pi)};
  \addplot[mspiB, domain=-1.6:2.6, samples=101] {exp(-x^2/2)/sqrt(2*pi)};
  \draw[black, dashed] (axis cs:0.5,0) -- (axis cs:0.5,1.55);
  \legend{{$f_{Z_n \mid H_1}$}, {$f_{Z_n \mid H_2}$}}
\nextgroupplot[title={$n = 16$: $P(e) = Q(2) \approx 0{,}023$}]
  \addplot[fill=mspiB!35, draw=none, domain=0.5:2.6, samples=60, forget plot] {4*exp(-8*x^2)/sqrt(2*pi)} \closedcycle;
  \addplot[fill=mspiA!35, draw=none, domain=-1.6:0.5, samples=60, forget plot] {4*exp(-8*(x-1)^2)/sqrt(2*pi)} \closedcycle;
  \addplot[mspiA, domain=-1.6:2.6, samples=161] {4*exp(-8*(x-1)^2)/sqrt(2*pi)};
  \addplot[mspiB, domain=-1.6:2.6, samples=161] {4*exp(-8*x^2)/sqrt(2*pi)};
  \draw[black, dashed] (axis cs:0.5,0) -- (axis cs:0.5,1.55);
  \legend{{$f_{Z_n \mid H_1}$}, {$f_{Z_n \mid H_2}$}}
\end{groupplot}
\end{tikzpicture}
\caption{Le densità della media campione $Z_n$ sotto le due ipotesi, con $\mu_1 = 1$, $\mu_2 = 0$ e $\sigma = 1$; le
code colorate oltre la soglia $\eta = \frac12$ sono gli errori. Con 16 misure le campane sono quattro volte più
strette (la deviazione standard è $\frac{\sigma}{\sqrt{n}}$) e si sovrappongono molto meno.}
\end{figure}

\gbox{In sintesi}{mspiA!80!black}{
\begin{itemize}
  \item Una \textbf{regola di decisione} $D(\vett{x}) \in \{1, \ldots, M\}$ è una partizione dei dati in regioni
        $\Omega_i$. Il \textbf{rischio bayesiano} $R = \sum_{i,j} C_{ij} P\{D = i, H_j\}$ pesa le decisioni con i costi;
        con costi 0 (giusto) e 1 (sbagliato) è la probabilità di errore.
  \item \textbf{Regola MAP} (minima $P(e)$): $L(\vett{x}) = \frac{p(\vett{x} \mid H_1)}{p(\vett{x} \mid H_2)}
        \decidi{H_1}{H_2} \frac{P(H_2)}{P(H_1)}$, cioè si decide per la massima probabilità a posteriori. Con ipotesi
        equiprobabili è la \textbf{regola ML}. Per dati continui si usano le densità.
  \item $P(e) = P(H_1) P(e \mid H_1) + P(H_2) P(e \mid H_2)$.
  \item Due sorgenti binarie: si confronta il \textbf{peso di Hamming} con una soglia; le probabilità di errore sono
        code di binomiali.
  \item Media gaussiana: si confronta la \textbf{media campione} con $\frac{\mu_1 + \mu_2}{2}$; $P(e) = Q\left(
        \sqrt{n} \frac{\mu_1 - \mu_2}{2\sigma} \right) \to 0$. Peso di Hamming e somma dei dati sono
        \textbf{statistiche sufficienti}.
\end{itemize}
}
